\part{Production, Maintenance, and Irreversibility}
\label{part:II}

Part II treats production as the transformation of conditions, capital as maintained capability, and irreversibility as a symmetric property of action and inaction.

\begin{housing}Roof and plumbing deteriorate; producing a habitable unit transforms the building's conditions; the cost of deferred maintenance is traced forward.\end{housing}

\chapter{Production as Transformation of Conditions}\label{ch:6}
Represents production through its outputs, consumed inputs, residuals, new obligations, and changes to future productive capacity. Compares a production function with a richer transition description and shows which questions each representation can answer.
\scaffold

\chapter{Capital as Maintained Capability}\label{ch:7}
Examines machines, buildings, knowledge, networks, and institutional arrangements as capabilities that require continued support. Capital valuation depends on complementary inputs, repairability, skills, and access, not on possession alone. A worked example compares two machines with identical output but different maintenance dependencies.
\scaffold

\chapter{Maintenance and the Reproduction of Capacity}\label{ch:8}
Develops maintenance schedules, deterioration processes, deferred work, and critical thresholds. A formal model shows how short-run expenditure reductions can reduce long-run feasible output when deterioration is nonlinear. Preventive maintenance is distinguished from restoration after failure.
\chspec{\Pref{P12}.}
\draftnote{Correction to carry over from earlier manuscripts: a nonpositive long-run average net accumulation rate does not ensure bounded debt, since $D(t)=\sqrt{1+t}-1$ is unbounded with average rate tending to zero. A bounded-debt theorem needs $\sup_{t\ge 0}\int_0^t(\rho_M-\rho_R)\,du<\infty$, or a reflected-backlog formulation.}
\scaffold

\chapter{Irreversible Decisions and the Value of Waiting}\label{ch:9}
Connects continuation economics to established work on irreversible investment and option value. Examines when waiting preserves useful information and when delay itself destroys opportunities. The result is a symmetric treatment of action and inaction as potentially irreversible.
\chspec{\Pref{P26}.}
\scaffold

\chapter{Residuals, Waste, and the Afterlives of Production}\label{ch:10}
Investigates demolition, abandoned infrastructure, pollution, obsolete software, and administrative residues. Asks what persists after an activity's advertised purpose has ended, and extends responsibility to the disposition of those residuals.
\scaffold
